\documentclass[11pt]{amsart}
\usepackage{amsmath,amssymb,amsthm,enumitem}
\newtheorem{theorem}{Theorem}
\newtheorem{conjecture}[theorem]{Conjecture}
\newtheorem{proposition}[theorem]{Proposition}
\newtheorem{remark}[theorem]{Remark}
\newcommand{\fg}{\mathfrak{g}}
\newcommand{\hbar}{\hslash}
\title{Chiral Coproduct for Non-Simply-Laced Types\\
via Quiver Folding}
\begin{document}
\maketitle

\section{The problem}

The chiral coproduct $\Delta\colon A \to A \otimes A$ on an
$E_1$-chiral quantum group is constructed for ADE types by
Jindal--Kaubrys--Latyntsev (arXiv:2603.21707) via critical CoHA
vertex coproducts.  For types $B_n$, $C_n$, $F_4$, $G_2$
(non-simply-laced), the construction is open.

\section{Yangian-level folding: complete}

$B_2 = \mathrm{fold}(A_3, \mathbb{Z}/2)$.  The involution
$\sigma(1) = 3$, $\sigma(2) = 2$ on the $A_3$ Dynkin diagram
induces $\sigma(T_{ij}(u)) = T_{5-i,5-j}(u)$ on the RTT generators
of $Y(\mathfrak{sl}_4)$.

\begin{proposition}[Yangian $\sigma$-equivariance]
The matrix coproduct on $Y(\mathfrak{sl}_4)$ is $\sigma$-equivariant:
$\Delta \circ \sigma = (\sigma \otimes \sigma) \circ \Delta$.
The fixed-point subalgebra
$Y(\mathfrak{sl}_4)^\sigma = Y(\mathfrak{so}_5)$ inherits the
matrix coproduct on the $5$-dimensional vector representation,
with $R$-matrix
$R^{B_2}(u) = 1 - \hbar P/u + \hbar Q/(u - 2)$
matching Molev and Jing--Liu--Molev.
\end{proposition}

\begin{proof}
$\Delta(\sigma(T_{ij})) = \Delta(T_{5-i,5-j})
= \sum_a T_{5-i,a} \otimes T_{a,5-j}
= \sum_b T_{5-i,5-b} \otimes T_{5-b,5-j}
= \sum_b \sigma(T_{ib}) \otimes \sigma(T_{bj})
= (\sigma \otimes \sigma)(\Delta(T_{ij}))$.
\end{proof}

\section{Chiral-level folding: two open steps}

The JKL vertex coproduct $\Delta^v$ on $V_k(\mathfrak{sl}_4)$
should restrict to $V_k(\mathfrak{so}_5) = V_k(\mathfrak{sl}_4)^\sigma$.

\textbf{Step 1: $\sigma$-equivariance of $\Delta^v$.}
The JKL construction uses the Hecke correspondence on
$\mathrm{Rep}(Q, \mathbf{d})$.  For the $A_3$ quiver with
$\mathbb{Z}/2$-involution, $\sigma$ acts on representation spaces
by permuting the vertex vector spaces.  The Hecke correspondence
(extension of representations) commutes with $\sigma$ by geometric
naturality.  The potential $W$ is $\sigma$-invariant (verified by
trace cyclicity).  Therefore $\Delta^v$ should commute with $\sigma$
at the level of the critical CoHA.

\textbf{Step 2: Factorisation through $V^\sigma \otimes V^\sigma$.}
$\sigma$-equivariance gives
$\Delta^v(V^\sigma) \subset (V \otimes V)^{\sigma \otimes \sigma}
= (V^+ \otimes V^+) \oplus (V^- \otimes V^-)$.
The $V^- \otimes V^-$ component must vanish for the coproduct to
factor through $V^\sigma \otimes V^\sigma$.  This requires
independent input beyond formal equivariance---it is the key
open step.

\section{The $G_2$ case}

$G_2 = \mathrm{fold}(D_4, \mathbb{Z}/3)$ (triality).  The cyclic
quotient singularity $\mathbb{C}^2/\mathbb{Z}_3$ does not admit a
crepant resolution, creating genuine combinatorial obstructions.
Third roots of unity appear in the spectral parameter.

\section{External input}

DeHority--Latyntsev (arXiv:2501.06643) develop the orthosymplectic
module theory of CoHAs with vertex coproducts in classical type.
Their twisted Yetter--Drinfeld module structure is compatible with
the folding framework and may provide the input needed for Step~2.

\end{document}
